\section{Literature Review}

The automated evaluation of academic papers represents a complex intersection of natural language processing, machine learning, and scholarly communication. This section surveys the evolution from traditional automated essay scoring (AES) systems to contemporary LLM-based approaches, examines recent advances in multi-agent consensus mechanisms, and identifies critical research gaps that motivate the design of HorizonBench.

\subsection{Automated Essay Scoring: From Feature Engineering to Deep Learning}

The field of Automated Essay Scoring traces its origins to the 1960s with Page's Project Essay Grade (PEG), which employed surface-level statistical features such as average sentence length, lexical complexity, and punctuation density to predict essay scores~\cite{page1966imminence}. While PEG demonstrated the feasibility of automated text assessment, its reliance on superficial linguistic proxies precluded meaningful evaluation of argumentative logic or content depth—limitations that become particularly pronounced when assessing academic papers where methodological rigor and intellectual contribution constitute primary value dimensions.

The Educational Testing Service's e-rater system represented a significant advancement, integrating grammar error detection, vocabulary usage analysis, and textual coherence features with Support Vector Machines and linear regression models~\cite{burstein1998automated}. Although e-rater achieved reasonable performance on standardized assessments such as TOEFL writing, its dependence on manually crafted features limited generalization capacity and rendered it inadequate for evaluating academic-specific dimensions including methodological soundness and experimental design validity.

The deep learning era brought substantial improvements in semantic understanding. Taghipour and Ng~\cite{taghipour2016neural} introduced a Bidirectional Long Short-Term Memory (BiLSTM) architecture capable of capturing sentence-to-document level semantic information, demonstrating notable improvements in coherence assessment for argumentative essays. Dong et al.~\cite{dong2017attention} further advanced the field by combining Convolutional Neural Networks with attention mechanisms to capture local linguistic patterns and supporting details. More recently, transformer-based approaches have achieved state-of-the-art results on benchmark datasets. Wang et al.~\cite{wang2022bert} proposed multi-scale BERT representations for AES, while hybrid approaches combining RoBERTa embeddings with handcrafted linguistic features have demonstrated QWK scores exceeding 0.92 on the ASAP dataset~\cite{iqbal2024hybrid}.

However, a comprehensive survey by Misgna et al.~\cite{misgna2024survey} emphasizes that despite these advances, critical aspects of writing evaluation—particularly feedback generation—have been largely overlooked. The survey identifies that while deep learning models excel at capturing general logical patterns, they cannot interpret specialized technical logic required for academic assessment. For instance, a model may detect that a passage lacks supporting evidence but cannot determine whether the design of a control group is appropriate or whether an ablation study sufficiently covers key variables—aspects central to academic evaluation.

\subsection{LLM-Based Automated Scholarly Paper Review}

The emergence of Large Language Models has catalyzed a paradigm shift toward what researchers term Automated Scholarly Paper Review (ASPR). A recent comprehensive survey~\cite{tan2025llm} documents this transformation, noting that LLMs' capacity for long-text multi-round conversation enables simulation of complete review processes, including author responses and reviewer interactions. The survey identifies supervised fine-tuning on peer review datasets as a promising direction, with systems like ReviewMT~\cite{tan2024reviewmt} converting the peer review process into multi-turn dialogues involving LLMs acting as authors, reviewers, and decision-makers.

Zhou et al.~\cite{zhou2024llm} conducted a comprehensive evaluation of LLM reliability as reviewers, introducing the RR-MCQ benchmark containing 197 review-revision multiple-choice questions from ICLR-2023. Their findings reveal a sobering reality: while LLMs can achieve passable decisions (>60\% accuracy) on single options, completely correct answers remain rare (approximately 20\%), with models demonstrating particular weakness in long paper processing, zero-shot scoring, and generating critical feedback comparable to human reviewers. This underscores the necessity for sophisticated prompting strategies and multi-model validation mechanisms.

The OpenReviewer system~\cite{openreviewer2024} represents a significant advance in specialized model development, fine-tuning Llama on 79,000 reviews from top ML conferences to generate human-like, high-quality reviews. Notably, the system demonstrates that specialized training on peer review data can overcome the tendency of general-purpose LLMs to generate overly positive assessments—a critical observation for system design. Similarly, the SEA framework~\cite{sea2024} introduces standardization across diverse academic venues, employing self-correction strategies to improve consistency between papers and generated reviews.

Recent work has also explored the simulation of peer review dynamics. The AGENTREVIEW framework~\cite{jin2024agentreview} at EMNLP 2024 integrates LLM agents to realistically simulate the peer review process, enabling controlled experiments to disentangle multiple variables affecting review outcomes. This framework demonstrates alignment with established sociological theories of peer review, validating the potential for agent-based simulation approaches.

\subsection{Multi-Agent Systems and Consensus Mechanisms}

The recognition that single-model approaches suffer from inherent biases and inconsistencies has motivated research into multi-agent consensus mechanisms. D'Arcy et al.~\cite{darcy2024marg} introduced MARG (Multi-Agent Review Generation), demonstrating that distributing paper text across specialized LLM agents engaging in internal discussion can substantially improve feedback quality. By employing a leader-worker-expert agent architecture, MARG reduced the rate of generic comments from 60\% to 29\% and doubled the generation of high-quality feedback compared to single-model baselines.

A comprehensive survey on LLM-based multi-agent systems~\cite{guo2024survey} identifies several critical mechanisms for effective collaboration: agent profiling with specialized expertise, structured communication protocols, and consensus-seeking strategies. The survey notes that rule-based strategies offer efficiency and predictability advantages by employing predefined rules governing agent interactions, ensuring straightforward implementation and facilitating debugging. Particularly relevant is the finding that peer review-inspired collaboration mechanisms, where agents critique and refine each other's outputs, can improve precision in complex reasoning tasks~\cite{xu2023peer}.

Recent work by Duan et al.~\cite{duan2024consensus} proposes integrating third-party LLMs to adjust attention weights through uncertainty estimation and confidence analysis, optimizing consensus formation in multi-agent systems. This approach demonstrates effectiveness in mitigating hallucinations that plague single-model deployments—a critical consideration for high-stakes academic evaluation applications.

The emergence of meta-review synthesis represents another important development. Traditional peer review employs Area Chairs who synthesize individual reviewer assessments into coherent decisions; analogously, computational approaches have begun exploring automated meta-review generation that resolves conflicts and identifies consensus points across multiple evaluations~\cite{li2023meta}.

\subsection{Benchmark Datasets and Evaluation Resources}

The development of appropriate benchmark datasets has proven essential for advancing automated review systems. The seminal PeerRead dataset~\cite{kang2018peerread} established the foundation, providing 14.7K paper drafts with accept/reject decisions from ACL, NIPS, and ICLR, alongside 10.7K expert peer reviews. Subsequent work has expanded this foundation: Ghosal et al.~\cite{ghosal2022peer} introduced Peer Review Analyze, a multi-layered dataset of 1,199 reviews manually annotated at the sentence level across four dimensions—Paper Section Correspondence, Paper Aspect Category, Review Functionality, and Review Significance.

The OpenReview platform has emerged as a critical resource, providing public access to submissions, reviews, author rebuttals, and final decisions for major conferences including ICLR, NeurIPS, and ICML~\cite{openreview}. Crucially, OpenReview data includes rejected submissions—information traditionally unavailable but essential for training systems to understand acceptance boundaries. The ORB dataset~\cite{orb2023} further systematizes this resource, aggregating over 36,000 submissions and 89,000 reviews from OpenReview and SciPost with standardized schemas.

For Retrieval-Augmented Generation evaluation, recent benchmarks have established rigorous assessment frameworks. The RAGBench dataset~\cite{ragbench2024} provides 100K examples with explainable labels across multiple dimensions including context relevance, answer faithfulness, and answer relevance. The ARES framework~\cite{ares2024} demonstrates that lightweight fine-tuned judges can accurately evaluate RAG systems while requiring only a few hundred human annotations, offering a promising direction for automated evaluation of retrieval-augmented academic assessment systems.

\subsection{Research Gaps and HorizonBench Positioning}

Despite substantial progress, our analysis identifies five critical gaps in current research that HorizonBench specifically addresses:

\textbf{1. Static vs. Dynamic Evaluation Baselines.} Existing systems evaluate papers against fixed, generic criteria without constructing domain-specific benchmarks. A paper in quantum computing cannot be meaningfully assessed against the same baseline as one in sentiment analysis. Current approaches either rely on broad rubrics that lack domain specificity or require manual construction of comparison sets—neither scalable nor comprehensive. HorizonBench addresses this through its Deep Research Agent (Phase 2), which dynamically constructs a Contextualized Benchmark by iteratively retrieving relevant prior work from Semantic Scholar, arXiv, and OpenReview, providing each paper with a tailored assessment baseline.

\textbf{2. Numerical Score Instability.} LLMs exhibit documented sensitivity to numerical outputs, producing different scores across repeated runs for identical inputs~\cite{zhou2024llm}. This instability fundamentally undermines reliability for assessment applications. While existing systems have attempted to mitigate this through prompt engineering or output constraints, none have fundamentally reconceptualized the scoring paradigm. HorizonBench adopts a grade-based rating system (Strong Accept through Strong Reject) aligned with real conference review practices, eliminating numerical instability while preserving assessment granularity.

\textbf{3. Single-Model Bias and Hallucination.} Current systems predominantly rely on single LLM instances, inheriting model-specific biases and hallucination tendencies. Although MARG~\cite{darcy2024marg} demonstrated benefits of multi-agent architectures, the agents employ the same underlying model (GPT-4), limiting diversity of perspectives. HorizonBench implements a truly heterogeneous multi-model consensus mechanism employing GPT, Claude, and DeepSeek as independent reviewers, with Gemini serving as an independent Meta Reviewer—a fourth model family ensuring no evaluator assesses its own outputs.

\textbf{4. Disconnect from Real Review Processes.} Existing systems generate assessments without modeling the structured deliberation that characterizes actual peer review. Real conferences employ Reviewers who provide independent assessments, followed by Area Chair synthesis that resolves conflicts and produces meta-reviews. Current automated systems collapse this multi-stage process into single-pass generation, losing the error-correction and deliberation benefits of the actual workflow. HorizonBench explicitly models this complete pipeline: independent parallel reviews (Phase 4a), conflict identification and resolution (Phase 4b), and synthesized meta-review with confidence estimation (Phase 4c).

\textbf{5. Limited Utilization of Rejection Data.} Perhaps most critically, existing systems predominantly train on accepted papers, lacking exposure to rejection rationales and boundary cases. The OpenReview platform provides unprecedented access to rejected submissions with detailed reviewer critiques—data essential for understanding what distinguishes acceptable from unacceptable work. HorizonBench integrates this resource directly into its few-shot prompting strategy, providing models with examples spanning Strong Accept through Strong Reject to calibrate assessment boundaries.

Furthermore, existing work has inadequately addressed the handling of novel or interdisciplinary research. When a system cannot find closely related prior work, traditional approaches simply fail or produce generic assessments. HorizonBench reconceptualizes this limitation as a feature: papers for which the Deep Research Agent finds sparse related literature represent potentially high-novelty contributions, triggering appropriate upward adjustment in innovation scores rather than assessment failure.

By addressing these gaps through an integrated four-phase architecture—Paper Understanding, Deep Research Agent, Multi-Dimensional Evaluation, and Multi-Model Consensus—HorizonBench establishes a comprehensive framework that more faithfully models, and potentially augments, the rigorous evaluation standards of top-tier academic venues.
